% Complete proof fragment FLH. Requires amsmath, amssymb, hyperref.
% Original amplitudes, full lattice, original metrics and tensor grades retained.
\section{Full-lattice null homology and the actual nilpotent receiver}
\label{sec:FLH}

\paragraph{The author source and the exact use made of it.}
Alain Connes and Caterina Consani,
\emph{Homological algebra in characteristic one},
\href{https://arxiv.org/abs/1703.02325v1}{arXiv:1703.02325v1},
is used through its original author file
\nolinkurl{homolalgpaper.tex}. The relevant source locators are
Proposition \texttt{Bmod1}, lines 610--640; Definition
\texttt{catb2}, lines 1607--1609; Definition
\texttt{ker\&coker}, lines 1757--1765; Proposition
\texttt{grandis} and its proof, lines 1862--1886; and the
definition of a null chain complex, lines 2680--2685.
The finite-subset quotient below is proved directly from their
kernel and cokernel universal properties. It does not invoke an
unread resolution theorem or identify the two coefficient categories.
Their involution is written \(\iota\) here because \(\sigma\)
already denotes the original support projection.

The programme inputs are the full coefficient laws HG1,
the synchronized maps OLR16--17, the actual rank-one operator
OLR13, and the complete tensor-ladder proofs OLM1--15 and
OLM26--28. All those objects and formulas used below are restated.
The new statements are the explicit localization, the two
different null-homology calculations, their exact reconstruction
maps, and their application to those original receivers.

\subsection{The original coefficient object and its actual null maps}

Let \(L\) be the full bounded distributive lattice with bottom
\(0_L\) and top \(1_L\). Do not replace it by a chain or by its
two endpoints. Put
\[
\begin{gathered}
S=G_L(\mathbb C)
 =\{(0,\lambda):\lambda\in L\}
       \cup\{(a,1_L):a\in\mathbb C\},\\
(a,\lambda)+(b,\mu)=(a+b,\lambda\vee\mu),\qquad
(a,\lambda)(b,\mu)=(ab,\lambda\wedge\mu),\\
\tau=(0,0_L),\quad z_\lambda=(0,\lambda),\quad
e_{\rm g}=z_{1_L},\quad \widehat a=(a,1_L),\quad
1_S=\widehat1,\quad \widehat0=e_{\rm g}.
\end{gathered}\tag{FLH1}
\]
These are the original component laws. They make \(S\) a
commutative semiring; closure follows because a nonzero
amplitude requires top support, and all distributive laws follow
from the ring and lattice laws. The symbol \(e_{\rm g}\) denotes
the original top-labelled zero, not the receiver vector \(e\).

For a complex vector space \(V\) retain the synchronized
\(S\)-semimodule
\[
\begin{gathered}
V_L^\uparrow
=\{(0,\lambda):\lambda\in L\}
       \cup\{(v,1_L):v\in V\},\\
(v,\lambda)+(w,\mu)=(v+w,\lambda\vee\mu),\qquad
(a,\eta)(v,\lambda)=(av,\eta\wedge\lambda),\\
\iota_V(v,\lambda)=(-v,\lambda),\qquad
\varepsilon_V(v,\lambda)=(0,\lambda)=e_{\rm g}(v,\lambda).
\end{gathered}\tag{FLH2}
\]
The semimodule laws follow from the same component computation.
Since the scalar field has characteristic zero,
\((V_L^\uparrow)^{\iota_V}=Z_L:=\{(0,\lambda):\lambda\in L\}\).
For every linear \(T:V\to W\) the original map
\[
T^\uparrow(v,\lambda)=(Tv,\lambda)
\tag{FLH3}
\]
is \(S\)-linear and involution-equivariant. Its composition law
is \((TU)^\uparrow=T^\uparrow U^\uparrow\), and its pointwise
sum law is \((T+U)^\uparrow=T^\uparrow+U^\uparrow\).
In particular \(0^\uparrow=\varepsilon_V\) with the appropriate
target. It is not the map constant at \(\tau\) when \(L\)
has a label different from \(0_L\).

There is a useful stronger statement for arbitrary
\(S\)-semimodules \(M\). The scalar \(u=\widehat{-1}\) defines
an involution, and
\[
M^\iota=\{m:e_{\rm g}m=m\},\qquad
f\text{ is null}\ \Longleftrightarrow\
e_{\rm g}f=f
\tag{FLH4}
\]
for \(S\)-linear \(f\), where null means that its image is fixed
by the involution. To prove the first equality, \(u e_{\rm g}
=e_{\rm g}\) gives one implication. Conversely \(um=m\)
gives \(e_{\rm g}m=(1_S+u)m=m+m=\widehat2m\).
Multiplication by the actual unit \(\widehat{1/2}\), and
\(\widehat{1/2}e_{\rm g}=e_{\rm g}\), give \(e_{\rm g}m=m\).
The second equality now follows pointwise. Composition with
linear maps preserves nullity, and every null map factors
through its target's fixed subsemimodule. This proves the
null ideal used here without assuming that \(S\) is Boolean.

\subsection{Support is both an idempotent reflection and a localization}

Give \(L\) its semiring operations \(\vee,\wedge\). The support
map and its section on labelled zeros are
\[
\sigma:S\longrightarrow L,\quad (a,\lambda)\longmapsto\lambda,
\qquad
j_L:L\longrightarrow S,\quad\lambda\longmapsto z_\lambda.
\tag{FLH5}
\]
The first map is a unital semiring homomorphism. The second
preserves addition and multiplication, but takes the unit to
\(e_{\rm g}\), so it is not a unital semiring section unless
\(e_{\rm g}=1_S\), which does not occur over \(\mathbb C\).
The actual identities are
\(\sigma j_L=I_L\) and \(j_L\sigma(s)=e_{\rm g}s\).

There is an exact unital localization
\[
 S[e_{\rm g}^{-1}]\ \simeq\ L,\qquad
 V_L^\uparrow[e_{\rm g}^{-1}]\ \simeq\ L.
\tag{FLH6}
\]
For the semiring assertion, let \(f:S\to A\) be unital and
suppose \(f(e_{\rm g})\) is invertible. It is idempotent, hence
equals \(1_A\). Thus
\(f(s)=f(e_{\rm g}s)=f(z_{\sigma(s)})\).
The unique factor is \(h(\lambda)=f(z_\lambda)\); it is a
unital semiring map because \(h(1_L)=1_A\), and the original
laws prove preservation of joins and meets. Conversely every
unital \(h:L\to A\) makes \(h\sigma(e_{\rm g})=1_A\).
This is the localization universal property, including the
case where inversion forces the zero semiring.

For the semimodule assertion, multiplication by \(e_{\rm g}\)
is the idempotent projection \(\varepsilon_V\) onto \(Z_L\).
Let \(A\) be an \(S\)-semimodule where \(e_{\rm g}\) acts
invertibly. It then acts as the identity. Every \(S\)-linear
\(f:V_L^\uparrow\to A\) satisfies
\(f(v,\lambda)=f(e_{\rm g}(v,\lambda))=f(0,\lambda)\).
It factors uniquely through
\(\sigma_V(v,\lambda)=\lambda\), with the \(L\)-action on
the target induced by \(\sigma\).
Conversely an \(L\)-linear factor gives such an \(S\)-linear
map. Under \(Z_L\simeq L\) its scalar action is exactly meet,
which proves the second universal property in FLH6.

The same map is the full additive idempotent reflection:
for every commutative idempotent monoid \(Q\), composition gives
\[
\operatorname{Hom}_{\rm join,0}(L,Q)
\ \simeq\
\operatorname{Hom}_{\rm add,0}(V_L^\uparrow,Q).
\tag{FLH7}
\]
Indeed, write \(a_v=f(v,1_L)\) and \(c=f(0,1_L)\).
The identities
\((v,1_L)+(0,1_L)=(v,1_L)\) and
\((v,1_L)+(-v,1_L)=(0,1_L)\) imply
\(a_v\vee c=a_v\) and \(a_v\vee a_{-v}=c\).
Thus \(c\le a_v\le c\), so \(a_v=c\).
The restriction \(h(\lambda)=f(0,\lambda)\) preserves joins
and bottom, and gives the unique factor in FLH7.
In particular an additive map into a Boolean module cannot
retain two different amplitudes at top support.

Every fibre and the full reconstruction are
\[
\begin{gathered}
\sigma_V^{-1}(\lambda)=\{(0,\lambda)\}\quad(\lambda\ne1_L),
\qquad
\sigma_V^{-1}(1_L)=V\times\{1_L\},\\
\pi_V^{-1}(v)=\{(v,1_L)\}\quad(v\ne0),\qquad
\pi_V^{-1}(0)=\{(0,\lambda):\lambda\in L\},\\
(\pi_V,\sigma_V):
 V_L^\uparrow\xrightarrow{\ \sim\ }
 \{(v,\lambda)\in V\times L:v=0\text{ or }\lambda=1_L\}.
\end{gathered}\tag{FLH8}
\]
The map in the last line and its inverse are literally the
ordered pair. Addition, scalar multiplication and involution
on its image are FLH2. Thus the amplitude/support joint object
recovers every state and operation, not merely the cardinalities
of the fibres. The same formulas apply to \(S\) with \(V=\mathbb C\).

\subsection{A literal finite complex and its full supported quotient}

Let
\[
0\longrightarrow C_m\xrightarrow{d_m}C_{m-1}
\longrightarrow\cdots\longrightarrow C_0\longrightarrow0,
\qquad d_nd_{n+1}=0
\tag{FLH9}
\]
be a finite complex of finite-dimensional complex vector spaces.
Finite here refers to its length and vector-space dimensions;
the underlying sets over \(\mathbb C\) are not finite objects
of the Boolean category.
Outside the displayed range put \(C_n=0\) and \(d_n=0\);
its synchronized completion has the full null object
\((0)_L^\uparrow=Z_L\) in those degrees. On every original
state, without removing its label,
\[
d_n^\uparrow d_{n+1}^\uparrow(v,\lambda)=(0,\lambda),
\qquad
\mathcal Z_{n,L}=(d_n^\uparrow)^{-1}(Z_L)
=(\ker d_n)_L^\uparrow .
\tag{FLH10}
\]
The composition is null by FLH4. The second formula is also
the null-kernel universal property in the entire category of
\(S\)-semimodules with the scalar involution: a map \(g\)
has \(d_n^\uparrow g\) null exactly when its image is contained
in this inverse image, so it factors uniquely through its
inclusion.

Write \(Z_n=\ker d_n\), \(B_n=\operatorname{im}d_{n+1}\subset Z_n\).
The full supported homology and quotient are
\[
\begin{gathered}
\mathcal H_{n,L}=(Z_n/B_n)_L^\uparrow,\qquad
q_n:\mathcal Z_{n,L}\longrightarrow\mathcal H_{n,L},\quad
(z,\lambda)\longmapsto([z],\lambda),\\
q_n^{-1}([z],1_L)=(z+B_n)\times\{1_L\},\qquad
q_n^{-1}(0,\lambda)=\{(0,\lambda)\}\quad(\lambda\ne1_L).
\end{gathered}\tag{FLH11}
\]
This is the null cokernel of the incoming differential viewed
as a map into \(\mathcal Z_{n,L}\). Here is its full universal
property, including targets not of synchronized vector form.
Let \(M\) be any \(S\)-semimodule and
\(g:\mathcal Z_{n,L}\to M\) be \(S\)-linear.
Then
\[
g\,d_{n+1}^\uparrow\text{ null}
\quad\Longleftrightarrow\quad
\exists!\,\bar g:\mathcal H_{n,L}\to M\text{ \(S\)-linear},
\quad g=\bar g q_n .
\tag{FLH12}
\]
To prove the forward implication, every \(b\in B_n\) has an
incoming preimage with top support. FLH4 gives
\(g(b,1_L)=e_{\rm g}g(b,1_L)=g(0,1_L)\).
Consequently
\[
g(z+b,1_L)=g(z,1_L)+g(b,1_L)
          =g(z,1_L)+g(0,1_L)=g(z,1_L).
\]
Thus \(g\) is constant on every fibre in FLH11, including
the singleton lower fibres, and defines the unique \(\bar g\).
Addition and scalar multiplication descend because \(q_n\)
is onto and preserves them. Conversely the image of
\(q_nd_{n+1}^\uparrow\) consists of labelled zero states,
which are fixed, and every equivariant map preserves fixed
states. This proves FLH12. Its kernel is \(B_{n,L}^\uparrow\),
so it also proves the exact normal boundary object here.

Every amplitude chain map \(f_n\) induces the exact supported
homology map
\[
([z],\lambda)\longmapsto([f_nz],\lambda).
\tag{FLH13}
\]
It is well-defined because it sends cycles to cycles and
boundaries to boundaries; FLH12 supplies the unique map.
These formulas preserve compositions and identities.
If \(C_n\) has its original positive Hermitian metric, the
amplitude quotient metric is, without altering that metric,
\[
\|[z]\|_{H_n}^{\,2}
 =\min_{b\in B_n}\|z+b\|_{C_n}^{\,2}
 =\|P_{Z_n\cap B_n^\perp}z\|_{C_n}^{\,2}.
\tag{FLH14}
\]
Finite-dimensional orthogonal decomposition of \(Z_n\) proves
existence, uniqueness of the minimizing representative, and
the equality. The supported pairing uses its original amplitude
pairing and the support meet; all labelled zeros have amplitude
norm zero but remain distinct.

\subsection{The precise Connes--Consani reflection and its homology}

Denote by \(\mathsf{BInv}\) the Connes--Consani category of
join-semilattices with bottom and a join-preserving involution.
Its null maps have fixed image. The object \(L\), with its
original joins and the identity involution, is a null object.
FLH6--7 give a functor from the actual synchronized category
to \(\mathsf{BInv}\):
\[
\mathcal R(V_L^\uparrow)=(L,I_L),\qquad
\mathcal R(T^\uparrow)=I_L,\qquad
\sigma_WT^\uparrow=\sigma_V .
\tag{FLH15}
\]
Equivalently it is localization followed by forgetting the
meet action. This is a genuine functor, since all identities
and compositions become the identity on the same \(L\).
In the reflected complex each differential is \(I_L\),
which is null because the involution is the identity.
Its null kernel is \(L\); the null cokernel of the incoming
identity is also \(L\), with quotient map the identity:
every equivariant map from this null object has fixed image
and hence already satisfies the cokernel condition. Thus
\[
H_n^{\rm CC}(\mathcal R(C^\uparrow))=(L,I_L)
 =\mathcal R(\mathcal H_{n,L})
\quad\text{in every degree}.
\tag{FLH16}
\]
This equality, with its explicit maps, is the homology
comparison for the additive Boolean reflection.
It is not an equality with the amplitude \(Z_n/B_n\).
Even if \(d_n=0\) on a nonzero vector space, all its amplitudes
are lost under FLH15, exactly as the fibres in FLH8 show.
One cannot instead apply the exchange-involution square
\(L\times L\) and keep the identity differential: its square
would be the nonnull identity whenever \(L\) has two distinct
elements. The null condition must be checked for the actual
involution and maps.

The faithful joint object retains the amplitude complex and
the reflected support complex together with the incidence
condition \(v=0\) or \(\lambda=1_L\). Its homology is exactly
\[
\{(h,\lambda)\in (Z_n/B_n)\times L:
                      h=0\text{ or }\lambda=1_L\},
\tag{FLH17}
\]
with all operations FLH2. This is proved by the quotient map
and every fibre in FLH11, not by interchanging an unspecified
limit with homology. The amplitude map on that quotient is
\(([z],\lambda)\mapsto[z]\), and its support map is the
identity on \(\lambda\). Therefore both the lost data and the
reconstruction are completely specified.

\subsection{An amplitude-retaining Boolean functor and its exact defect}

There is also a different, nonadditive-on-states, functor to
\(\mathsf{BInv}\) that retains each original state. Define
\[
\begin{gathered}
\mathcal P_L(V)=
\{(A,\lambda): A\subset V\text{ finite},\ 0\in A,\
                     A\ne\{0\}\Rightarrow\lambda=1_L\},\\
(A,\lambda)\vee(A',\mu)=(A\cup A',\lambda\vee\mu),\qquad
0_{\mathcal P}=(\{0\},0_L),\qquad
\iota(A,\lambda)=(-A,\lambda),\\
\mathcal P_L(T)(A,\lambda)=(T(A),\lambda),\qquad
j_V(v,\lambda)=(\{0,v\},\lambda).
\end{gathered}\tag{FLH18}
\]
Union and join prove every Boolean-module law; negation is
an involution preserving those operations. A linear \(T\)
sends \(0\) to \(0\), preserves unions under direct image,
and commutes with negation. It therefore gives a morphism in
\(\mathsf{BInv}\), with compositions and identities preserved.
The map \(j_V\) is injective, pointed, equivariant and natural:
\(\mathcal P_L(T)j_V=j_WT^\uparrow\). It does not preserve
vector addition: for a nonzero \(v\),
\(\{0,v\}\cup\{0,-v\}\ne\{0\}\).
Thus FLH18 asserts no additive identification forbidden by FLH7.
No scalar-distributive \(S\)-module structure on finite sets
is being assumed.

The null complex and its cycles are now calculated exactly:
\[
\begin{gathered}
\mathcal P_L(d_n)\mathcal P_L(d_{n+1})(A,\lambda)
   =(\{0\},\lambda),\\
\mathcal K_n=
\{(A,\lambda)\in\mathcal P_L(C_n):
                              d_n(A)=-d_n(A)\}.
\end{gathered}\tag{FLH19}
\]
The first image is fixed. The second line is the literal
inverse image of fixed points, hence the Connes--Consani
null kernel and has their stated universal property.
Unions of such sets still satisfy the displayed condition,
so it is a subobject.

Here is an evaluated normal form for its entire null homology.
Let \(\mathscr O(B_n)=(B_n\setminus\{0\})/\{b\sim-b\}\);
its elements are exact two-element sign orbits, not projective
lines and not scalar orbits. Put
\[
\begin{gathered}
\mathcal Q_n=\{(A_{\rm out},\Omega,\lambda):
 A_{\rm out}\subset C_n\setminus B_n\text{ finite},\
 \Omega\subset\mathscr O(B_n)\text{ finite},\\
\{0\}\cup d_n(A_{\rm out})
        =-\bigl(\{0\}\cup d_n(A_{\rm out})\bigr),\
(A_{\rm out}\ne\varnothing\text{ or }\Omega\ne\varnothing)
                      \Rightarrow\lambda=1_L\},\\
q_n^{\rm CC}(A,\lambda)
 =\bigl(A\setminus B_n,\
       \{\{b,-b\}:b\in A\cap(B_n\setminus\{0\})\},\
       \lambda\bigr).
\end{gathered}\tag{FLH20}
\]
The joins are union in the first two coordinates and join
in the last. Its involution sends
\((A_{\rm out},\Omega,\lambda)\) to
\((-A_{\rm out},\Omega,\lambda)\).
These formulas define a Boolean module with involution.
The map \(q_n^{\rm CC}\) preserves its joins, bottom and
involution. It is onto: choose one representative of each
of the finitely many orbits in \(\Omega\), append \(A_{\rm out}\)
and \(0\), and use that \(d_n(B_n)=0\).

Its full fibre equivalence and universal property are
\[
\begin{gathered}
(A,\lambda)\sim_{\rm CC}(A',\mu)
\ \Longleftrightarrow\
\lambda=\mu,\quad A\setminus B_n=A'\setminus B_n,\quad
\mathscr O(A\cap B_n)=\mathscr O(A'\cap B_n),\\
H_n^{\rm CC}(\mathcal P_L(C))\simeq\mathcal Q_n.
\end{gathered}\tag{FLH21}
\]
Here the orbit expression omits zero as in FLH20.
We prove the cokernel property rather than infer it from ranks.
The incoming image in \(\mathcal K_n\) is precisely
\(\mathcal P_L(B_n)\): lift its finitely many vectors
individually through \(d_{n+1}\), preserving the label,
and retain \(0\). Its image under \(q_n^{\rm CC}\) has empty
outside set and therefore is fixed.

Let \(f:\mathcal K_n\to M\) be any morphism in
\(\mathsf{BInv}\) for which the incoming image is null.
For every nonzero \(b\in B_n\), equivariance and nullity give
\[
f(\{0,b\},1_L)=f(\{0,-b\},1_L).
\tag{FLH22}
\]
Joining an arbitrary element of \(\mathcal K_n\) to both
sides permits replacing \(b\) by \(-b\), or adding the other
sign and then removing a redundant sign. Removing boundary
elements from a cycle set leaves a cycle set because their
differential is zero. Consequently any two finite sets in
one fibre of FLH20 are connected by these finitely many
replacements, keeping their outside sets and support fixed.
The all-zero set with a lower label is a singleton fibre.
Thus \(f\) is constant on every fibre and factors uniquely
through the surjection \(q_n^{\rm CC}\). The resulting map
preserves joins and the involution by surjectivity.
Conversely every factor makes the incoming image fixed.
This proves the cokernel universal property for all targets
of the Connes--Consani category, and proves FLH21.
In particular a nonzero boundary sign orbit becomes fixed;
it is not deleted or identified with the labelled empty orbit.
Every fibre is more explicitly
\[
\left\{\left(\{0\}\cup A_{\rm out}
                   \cup\bigcup_{\omega\in\Omega}F_\omega,\lambda\right):
             \varnothing\ne F_\omega\subseteq\omega
                             \text{ for each }\omega\in\Omega\right\}.
\]
The sign orbits are disjoint two-element sets, so the fibre has
exactly \(3^{|\Omega|}\) elements. In particular no amplitude
translation by a boundary has entered this quotient.
There is a canonical section in the Boolean category:
\[
c_n(A_{\rm out},\Omega,\lambda)
 =(\{0\}\cup A_{\rm out}\cup\bigcup_{\omega\in\Omega}\omega,\lambda).
\]
It preserves unions, labels and involution, and
\(q_n^{\rm CC}c_n=I\).
Its other composition is the exact idempotent
\((A,\lambda)\mapsto(A\cup[-(A\cap B_n)],\lambda)\).
It adds the missing opposite boundary signs and changes
no other vector.
Consequently the fixed subobject of \(\mathcal Q_n\) consists
precisely of the triples with \(A_{\rm out}=-A_{\rm out}\).
The normal image of the incoming map, namely the null kernel
of its null cokernel, is precisely the subset of
\(\mathcal K_n\) with symmetric outside part. It equals
\(\mathcal P_L(B_n)\vee\mathcal K_n^\iota\):
one inclusion follows by taking outside parts of a union;
for the other decompose a cycle set with symmetric outside
part into \(A\cap B_n\) and \(\{0\}\cup(A\setminus B_n)\),
the latter being fixed. A nontrivial such union has top
label, and the zero sets already lie in the incoming image
at every label. This proves the asserted equality with
no label omitted.
The construction is natural under every original chain map.
For \(f_n:C_n\to C'_n\), take the canonical representative
\(\{0\}\cup A_{\rm out}\cup\bigcup\Omega\), apply \(f_n\)
to every vector, and take its FLH20 normal form for \(B'_n\).
The chain identity preserves the balanced-image condition
and sends \(B_n\) into \(B'_n\); opposite boundary representatives
therefore still give the same orbit, or zero, in the target.
This proves well-definition and gives the actual induced
map on the complete normal-form homology.

This object contains a precisely describable excess over
additive cycles. The original cycles are recovered by the
exact marked-state identity
\[
j_{C_n}((Z_n)_L^\uparrow)
 =j_{C_n}((C_n)_L^\uparrow)\cap\mathcal K_n.
\tag{FLH23}
\]
Indeed \(\{0,d_nv\}\) is sign-invariant over \(\mathbb C\)
if and only if \(d_nv=0\). The original boundary marked states
are \(\{j_{C_n}(b,1_L):b\in B_n\}\) together with all labels.
Their additive quotient must use the original addition graph:
\((z,\lambda)\) and \((z',\mu)\) represent the same supported
class exactly when \(\lambda=\mu\), and either that label is
not top, or \(z'-z\in B_n\). This is precisely FLH11--12.
The singleton marking, its original addition and scalar
graphs, and its support therefore reconstruct the additive
homology from the actual complex. They are additional
specified structure; they are not recovered by declaring
the Boolean join to be vector addition.

For example, in a two-term complex \(C_1\xrightarrow{T}C_0\),
choose \(Tv\ne0\) and \(0\ne w\in\ker T\). Then
\[
A=\{0,v,-v+w\},\qquad T(A)=\{0,Tv,-Tv\}.
\tag{FLH24}
\]
The set \(A\) is a Boolean null cycle. It is not fixed by
negation: \(-v\) is none of \(0,v,-v+w\), since characteristic
zero, \(Tv\ne0\), and \(w\ne0\) exclude the three equalities.
There is no incoming amplitude boundary in degree one, so
FLH21 keeps this nonfixed class. Neither nonzero vector in
\(A\) is an additive cycle. Thus this functor supplies a
genuine characteristic-one null complex but does not silently
identify its homology with FLH11. Its extra objects and the
maps recovering the actual cycles are now explicit.

\subsection{The rank-one receiver, with its original metric}

Keep the original receiver \(E\), positive metric \(G_N\),
unit vectors \(e,f\), and coefficient \(\epsilon>0\):
\[
\begin{gathered}
\langle v,w\rangle_N=v^*G_Nw,\quad
\|e\|_N=\|f\|_N=1,\quad e^\dagger f=0,\quad
R=\epsilon f e^\dagger,\quad R^2=0,\\
W=\mathbb Ce\mathbin\perp\mathbb Cf,\quad
U=W^\perp,\quad E=W\mathbin\perp U,\quad \dim E=q .
\end{gathered}\tag{FLH25}
\]
In the literal two-term complex \(C_1=E\xrightarrow{R}C_0=E\),
completed outside these degrees by the supported null object,
the complete supported homology is
\[
\mathcal H_{1,L}=(\mathbb Cf\mathbin\perp U)_L^\uparrow,\qquad
\mathcal H_{0,L}=(E/\mathbb Cf)_L^\uparrow
 \simeq(\mathbb Ce\mathbin\perp U)_L^\uparrow .
\tag{FLH26}
\]
The last isomorphism sends a coset to its original orthogonal
representative \(v-ff^\dagger v\), keeping its label.
Indeed \(Rv=\epsilon f(e^\dagger v)\) has kernel \(e^\perp\)
and image \(\mathbb Cf\); FLH11 proves the supported quotient
and FLH14 proves its metric. On \(a e+b f+u\) the quotient
norm is \(|a|^2+\|u\|_N^2\), whereas the kernel norm on
\(bf+u\) is \(|b|^2+\|u\|_N^2\).
The actual inverse from the image to the orthogonal source
complement is \(bf\mapsto (b/\epsilon)e\), with squared norm
\(|b|^2/\epsilon^2\).

Here a three-term complex is also genuine:
\(E\xrightarrow{R}E\xrightarrow{R}E\).
Its middle homology is \(U_L^\uparrow\), its top homology is
\((\mathbb Cf\perp U)_L^\uparrow\), and its bottom homology is
\((\mathbb Ce\perp U)_L^\uparrow\), with the same quotient maps.
This follows from \(R^2=0\), the displayed kernel and image,
and their original orthogonal decomposition.
The scalar \(\epsilon\) remains in its differential and inverse,
although it does not change those subspaces.

The Boolean homology of the two-term complex is given
without an unidentified quotient by
\[
\begin{aligned}
H_1^{\rm CC}(\mathcal P_L(C))
 &=\{(A,\lambda):A\subset E\text{ finite},\ 0\in A,\
       \epsilon f\,e^\dagger A=-\epsilon f\,e^\dagger A\},\\
H_0^{\rm CC}(\mathcal P_L(C))
 &=\{(A_{\rm out},\Omega,\lambda):
 A_{\rm out}\subset E\setminus\mathbb Cf\text{ finite},\
 \Omega\subset(\mathbb Cf\setminus\{0\})/\{\pm1\}\text{ finite}\}.
\end{aligned}\tag{FLH27}
\]
In both lines retain exactly the support rule of FLH18 or
FLH20, respectively. Their operations and involutions are
the displayed normal-form operations, not vector operations.
Taking \(v=e\) and \(0\ne w\in U\) gives the explicit nonfixed
extra cycle FLH24 in the actual receiver whenever \(q>2\),
as it is in the programme.
The additive reflection of both homologies in FLH26 remains
the same null \(L\), by FLH16.

\subsection{The mixed tensor receiver as a genuine two-term complex}

Retain \(k\ge17\), \(k\equiv1\pmod4\), \(q=(k+1)^2\), the
original cutoff \(q-1\le N\le2q\), and \(G_N\) of OLM1.
For \(t\ge1\) the unchanged operator and metric are
\[
E_t=E^{\otimes t},\quad G_t=G_N^{\otimes t},\qquad
X=R^{[t]}=\sum_{a=1}^t
 I^{\otimes(a-1)}\otimes R\otimes I^{\otimes(t-a)} .
\tag{FLH28}
\]
For \(t\ge2\), this operator is not a differential repeated
in adjacent degrees:
\[
X^2(e^{\otimes t})
=2\epsilon^2\sum_{1\le a<b\le t}
 e\otimes\cdots\otimes f_a\otimes\cdots\otimes f_b
                       \otimes\cdots\otimes e\ne0 .
\tag{FLH29}
\]
The distinct tensor words in the sum are orthogonal in the
original metric, each coefficient is \(2\epsilon^2\ne0\),
and hence the claimed nonvanishing follows. We use the
literal two-term complex \(C_1=E_t\xrightarrow{X}C_0=E_t\).
The next differential is the zero amplitude map with its
full support lift, so this is a valid null complex for every
\(t\). No sign change is made to the slot-sum operator.

For every actual active-slot set \(J\subseteq\{1,\ldots,t\}\),
put \(j=|J|\). The original orthogonal sector is
\(E_J=W^{\otimes j}\otimes U^{\otimes(t-j)}\), embedded in its
actual slots in their increasing order. Let \(V_{j,r}\) be
the span of words with \(r\) copies of \(f\), and \(X_j\) the
sum of \(R\) over these \(j\) slots. With \(Y_j=X_j^\dagger\),
the complete primitive spaces and factors are
\[
\begin{gathered}
\mathcal P_{j,r}=\ker(Y_j:V_{j,r}\to V_{j,r-1}),\quad
0\le r\le\lfloor j/2\rfloor,\quad d=j-2r,\\
\dim\mathcal P_{j,r}=\binom jr-\binom j{r-1},\quad
s_{d,a}=\epsilon^{2a}a!\frac{d!}{(d-a)!},\\
E_t=\bigoplus_{J,r}\ \bigoplus_{a=0}^{d}{}^\perp
 X_j^a\mathcal P_{j,r}\otimes U^{\otimes(t-j)},\qquad
\langle X_j^au,X_j^bv\rangle
 =\mathbf1_{a=b}s_{d,a}\langle u,v\rangle .
\end{gathered}\tag{FLH30}
\]
The sector embeddings keep the original order of all slots.
For completeness, the one-slot commutator gives
\([Y_j,X_j]|_{V_{j,s}}=\epsilon^2(j-2s)I\).
It implies injectivity below the middle grade and, by the
slotwise \(e,f\) exchange, surjectivity above it. Rank-nullity
gives the primitive dimension. Commuting \(Y_j\) through
\(X_j^a\) on a primitive vector yields
\(Y_jX_j^au=\epsilon^2a(d-a+1)X_j^{a-1}u\).
Induction gives the full factor \(s_{d,a}\), and the next
power has zero norm, hence vanishes. Different grades and
different primitive ladders are orthogonal: moving the smaller
number of powers to the other entry leaves either a positive
power paired with a primitive vector or different grades.
At total grade \(s\), their dimensions telescope as
\(\sum_{r=0}^{\min(s,j-s)}[\binom jr-\binom j{r-1}]
=\binom js\). Thus they exhaust every grade and every slot
sector. This proves the decomposition and metric in FLH30
with the unchanged vectors and coefficients.

The full homology is now evaluated as actual subspaces and
quotient metrics:
\[
\begin{gathered}
H_1(C)=\ker X
 =\bigoplus_{J,r}{}^\perp
       X_j^d\mathcal P_{j,r}\otimes U^{\otimes(t-j)},\\
H_0(C)=E_t/\operatorname{im}X
 \xrightarrow{\ \sim\ }
 \bigoplus_{J,r}{}^\perp
       \mathcal P_{j,r}\otimes U^{\otimes(t-j)},\\
\mathcal H_{n,L}=H_n(C)_L^\uparrow\quad(n=0,1),\\
\left\|\sum_{J,r}X_j^d u_{J,r}\right\|_{H_1}^{\,2}
   =\sum_{J,r}\epsilon^{2d}(d!)^2\|u_{J,r}\|^2,\qquad
\|[v]\|_{H_0}^{\,2}=\sum_{J,r}\|u_{J,r,0}\|^2 .
\end{gathered}\tag{FLH31}
\]
Here \(v=\sum_{J,r,a}X_j^a u_{J,r,a}\) is its actual full
decomposition, and the second line is the original orthogonal
representative of the coset. On each ladder \(X\) sends level
\(a\) to \(a+1\) and kills exactly its top level \(d\).
Its image is exactly the sum of levels \(1,\ldots,d\).
Independence and the metric in FLH30 prove every equality
and the quotient-minimum assertion in FLH31.

All its ranks, including every inactive-slot multiplicity,
are therefore
\[
\begin{gathered}
b_t:=\dim H_1(C)=\dim H_0(C)
 =\sum_{j=0}^t\binom tj(q-2)^{t-j}
                         \binom j{\lfloor j/2\rfloor},\\
\operatorname{rank}X=q^t-b_t,\qquad
M_{t;j,r}=\binom tj(q-2)^{t-j}
                         \left(\binom jr-\binom j{r-1}\right).
\end{gathered}\tag{FLH32}
\]
Each primitive vector contributes one top kernel vector and
one bottom quotient vector. Summing its exact dimension
over \(r\) telescopes to the middle binomial coefficient.
There are \(\binom tj\) actual choices of \(J\) and
\((q-2)^{t-j}\) inactive-factor dimensions. This proves FLH32,
including \(j=0,d=0\).
The support fibre over every non-top label in FLH31 remains
one labelled zero; its top fibre is the entire displayed
homology space, even if a particular summand has dimension zero.

The actual inverse on the image and all its metric factors
are also explicit:
\[
\begin{gathered}
X^+\!\left(\sum_{J,r}\sum_{b=1}^{d}X_j^b u_{J,r,b}\right)
 =\sum_{J,r}\sum_{b=1}^{d}X_j^{b-1}u_{J,r,b},\\
\|X^+v\|^2=\sum_{J,r}\sum_{b=1}^{d}
                              s_{d,b-1}\|u_{J,r,b}\|^2,\qquad
\frac{s_{d,b}}{s_{d,b-1}}=\epsilon^2b(d-b+1).
\end{gathered}\tag{FLH33}
\]
This inverse lands in the source orthogonal complement of
\(\ker X\), so it is the attained minimum lift of \(v\).
Orthogonality in FLH30 proves uniqueness and its norm.
The superscript \(+\) denotes this original-metric
Moore--Penrose inverse; it does not divide away \(\epsilon\)
from the original differential.
All maps and inverses in FLH31--33 lift by FLH3 and retain
their label.

Finally the complete Boolean null homology of this same
two-term tensor complex, with no tensor-square-zero error, is
\[
\begin{aligned}
H_1^{\rm CC}(\mathcal P_L(C))
 &=\{(A,\lambda)\in\mathcal P_L(E_t):X(A)=-X(A)\},\\
H_0^{\rm CC}(\mathcal P_L(C))
 &=\{(A_{\rm out},\Omega,\lambda):
 A_{\rm out}\subset E_t\setminus\operatorname{im}X\text{ finite},\
 \Omega\subset(\operatorname{im}X\setminus\{0\})/\{\pm1\}
                          \text{ finite}\}.
\end{aligned}\tag{FLH34}
\]
Retain the support rules and operations of FLH20.
The image appearing here is exactly the sum of levels
\(1,\ldots,d\) in every sector of FLH30. Thus every set,
sign orbit, degree, kernel and quotient in FLH34 is determined
by the original operator. FLH23 selects its original additive
cycles, and FLH11 then gives the supported additive quotient
with all metrics FLH31. Localization applies to the
synchronized additive complex and its supported homology,
and sends those objects to the null \(L\) in FLH16.
The Boolean normal-form object instead has the explicitly
defined support morphism
\(\rho_{\rm CC}(A_{\rm out},\Omega,\lambda)=\lambda\)
to \((L,I_L)\). It preserves joins and bottom and commutes
with involution by FLH20. Its fibre at a non-top label is
the single \((\varnothing,\varnothing,\lambda)\), and its
top fibre is the entire displayed set of allowed first two
coordinates at that label. No \(S\)-semimodule localization
of this Boolean object is asserted.

The first line of FLH34 has a complete description by actual
affine fibres, beyond its kernel presentation. Put \(K=\ker X\)
as explicitly decomposed in FLH31, and let \(X^+\) be FLH33.
Every one of its elements is uniquely of the form
\[
\begin{gathered}
A=A_0\ \cup\!
       \bigcup_{w\in T}\bigl(X^+w+U_w\bigr),\qquad
T\subset\operatorname{im}X\setminus\{0\}\text{ finite},\
T=-T,\\
A_0\subset K\text{ finite},\quad0\in A_0,\qquad
\varnothing\ne U_w\subset K\text{ finite}\quad(w\in T),\\
(A_0\ne\{0\}\text{ or }T\ne\varnothing)
                  \Rightarrow\lambda=1_L,\\
\iota:(A_0,T,(U_w),\lambda)
       \longmapsto(-A_0,T,(-U_{-w})_{w\in T},\lambda).
\end{gathered}\tag{FLH35}
\]
Indeed the minimum inverse satisfies \(XX^+w=w\), so
the complete vector fibre is \(X^{-1}(w)=X^+w+K\).
The displayed union has image \(\{0\}\cup T\), and hence
is a null cycle. Conversely, for a cycle \(A\), take
\(T=X(A)\setminus\{0\}\), \(A_0=A\cap K\), and
\(U_w=(A\cap X^{-1}(w))-X^+w\).
These sets have the displayed properties and recover \(A\).
The fibres of different \(w\) are disjoint, which proves
uniqueness. Linearity of \(X^+\) proves the involution formula.
In the rank-one case the literal fibre section is
\(X^+(bf)=(b/\epsilon)e\); in every tensor sector it is
the unaltered level-lowering map FLH33.
This retains the numerical coefficients in all the Boolean
cycle fibres. In the second line of FLH34, the outside vectors
are precisely those with at least one nonzero bottom coordinate
\(u_{J,r,0}\), whereas the boundary sign orbits are the nonzero
vectors whose every bottom coordinate is zero.
These statements follow from the evaluated image in FLH31.
The Boolean category alone assigns no Hermitian norm to
finite sets; all original vector norms remain the exact
ones in FLH30--33 and in the joint marked-state data.

\subsection{The joint diagram and the original free masks}

\begin{figure}[ht]
\centering
\[
\begin{array}{ccccc}
(Z_n)_L^\uparrow &\xrightarrow{\ q_n\ }&
(Z_n/B_n)_L^\uparrow&\xrightarrow{\ \sigma_H\ }&(L,I_L)\\
\pi_Z\downarrow&&\pi_H\downarrow&&\\
Z_n&\xrightarrow{\ z\mapsto[z]\ }&Z_n/B_n&&
\end{array}
\]
\[
\begin{array}{ccccc}
(Z_n)_L^\uparrow&\xrightarrow{\ j_Z\ }&
\mathcal K_n&\xrightarrow{\ q_n^{\rm CC}\ }&
\mathcal Q_n\\
(z,\lambda)&\longmapsto&
(\{0,z\},\lambda)&& (A_{\rm out},\Omega,\lambda)
\end{array}
\]
\caption{The two proved quotients of the original finite complex.
The upper row retains the amplitude coset and every support label
(FLH11--17). The bottom row is the Connes--Consani Boolean null
quotient, retaining outside-boundary sets and boundary sign orbits
(FLH18--23); it is not an arrow identifying those sign orbits with
amplitude cosets. The map \(j_Z\) marks the original individual
states; it is pointed and equivariant, and is not additive.
Connes--Consani's
kernel and cokernel definitions are at the author-source locators
given at the start of this section.}
\label{fig:FLH-joint}
\end{figure}

The square of the first two columns commutes by FLH11.
The support on the upper quotient is reconstructed jointly
with the amplitude, as in FLH17; the Boolean map in the last
row is the different universal quotient proved in FLH21--22.

In the original positive-dimensional free coordinate module,
the map to a synchronized state and its whole fibre are
\[
\begin{gathered}
P_{\rm syn}(v,\boldsymbol\mu)
       =(v,\bigvee_i\mu_i),\\
P_{\rm syn}^{-1}(v,\lambda)
 =\{(v,\boldsymbol\mu):
       \bigvee_i\mu_i=\lambda,\quad
       v_i\ne0\Rightarrow\mu_i=1_L\}.
\end{gathered}\tag{FLH36}
\]
The join is unchanged when synchronization is repeated,
so this is the original idempotent correspondence OLR17.
Combining it with FLH10--11 gives every free mask over every
cycle and supported homology fibre: the amplitude belongs
to the stated cycle or affine boundary coset, while the masks
satisfy exactly FLH36. No change of free-semimodule basis
to a ladder basis is claimed. The ladder maps are applied
after synchronization and have their proved amplitude inverse.
If the entire vector space has dimension zero, the zero-coordinate
free module has one state, whereas \((0)_L^\uparrow\) has
all labels; FLH36 is not declared onto in that case.
The actual receiver and its positive tensor powers have
dimensions \(q\) and \(q^t>0\), so their full masks are covered.

The source localization, the Connes--Consani null quotient,
the additive supported quotient, and the actual nilpotent
receiver have thus been related by explicit maps and complete
fibres. None of FLH1--36 supplies a Frobenius purity assertion,
an RH contradiction, or a counterexample to RH. The calculated
amplitude loss under localization and the additional Boolean
cycles are part of the resulting objects, with the reconstruction
and receiver calculations above retaining their full data.
